\subsection{模型建立}
我们在问题四中沿用问题三的方程组和参数集，仅将判别工具替换为 Benettin 型有限时间最大李雅普诺夫指数 $\lambda_{\max}$。数值实现中，先由 Hastings-Powell 方程写出解析 Jacobian，再将状态方程与切向方程用四阶 Runge--Kutta 同步积分，切向向量在每个时间步进行一次归一化；“瞬态截断 40”指前 40 个时间单位的轨道不进入指数累计。基准设置取初值 $(0.8,0.2,0.2)$、步长 $dt=0.01$ 和总积分时长 90；同时将总积分时长扩展到 150 和 200，并在 $K=3.20$ 与 $K=3.22$ 附近进行局部步长和初值复核，用以检验正指数起点对数值设置的敏感性。

